% ============================================================
% 初升高衔接讲义 — 不等式与二次不等式
% 编译命令：xelatex -interaction=nonstopmode 衔接3_不等式与二次不等式.tex
% ============================================================
\documentclass[12pt, a4paper]{ctexart}
\usepackage[margin=2.2cm]{geometry}
\usepackage{amsmath, amssymb}
\usepackage{enumitem}
\usepackage{tcolorbox}
\usepackage{booktabs}
\usepackage{array}
\usepackage{fancyhdr}
\usepackage{xcolor}
\usepackage{tikz}
\usepackage{hyperref}
\usepackage{graphicx}
\usepackage{multirow}
\usepackage{pgfplots}
\pgfplotsset{compat=1.18}

% 颜色定义
\definecolor{mainblue}{RGB}{41, 98, 168}
\definecolor{lightblue}{RGB}{230, 240, 255}
\definecolor{tipgreen}{RGB}{34, 139, 87}
\definecolor{lightgreen}{RGB}{235, 250, 240}
\definecolor{warnred}{RGB}{200, 50, 50}
\definecolor{lightred}{RGB}{255, 240, 240}
\definecolor{examorange}{RGB}{210, 120, 20}
\definecolor{lightorange}{RGB}{255, 248, 235}
\definecolor{apppurple}{RGB}{120, 60, 170}
\definecolor{lightpurple}{RGB}{245, 238, 255}

% 页眉页脚
\pagestyle{fancy}
\fancyhf{}
\fancyhead[L]{\small\textcolor{mainblue}{\textbf{初升高衔接 · 不等式与二次不等式}}}
\fancyhead[R]{\small\textcolor{mainblue}{衔接讲义三}}
\fancyfoot[C]{\small —\ \thepage\ —}
\renewcommand{\headrulewidth}{0.6pt}
\renewcommand{\headrule}{\hbox to\headwidth{\color{mainblue}\leaders\hrule height \headrulewidth\hfill}}
\renewcommand{\footrulewidth}{0pt}

% tcolorbox 样式
\newtcolorbox{knowledgebox}[1][]{%
  colback=lightblue, colframe=mainblue,
  fonttitle=\bfseries\large, title=#1,
  arc=3mm, boxrule=0.8pt,
  left=8pt, right=8pt, top=6pt, bottom=6pt
}

\newtcolorbox{tipbox}[1][]{%
  colback=lightgreen, colframe=tipgreen,
  fonttitle=\bfseries, title=#1,
  arc=2mm, boxrule=0.6pt,
  left=6pt, right=6pt, top=4pt, bottom=4pt
}

\newtcolorbox{warnbox}[1][]{%
  colback=lightred, colframe=warnred,
  fonttitle=\bfseries, title=#1,
  arc=2mm, boxrule=0.6pt,
  left=6pt, right=6pt, top=4pt, bottom=4pt
}

\newtcolorbox{examplebox}[1][]{%
  colback=lightorange, colframe=examorange,
  fonttitle=\bfseries, title=#1,
  arc=3mm, boxrule=0.8pt,
  left=8pt, right=8pt, top=6pt, bottom=6pt
}

\newtcolorbox{appbox}[1][]{%
  colback=lightpurple, colframe=apppurple,
  fonttitle=\bfseries, title=#1,
  arc=3mm, boxrule=0.8pt,
  left=8pt, right=8pt, top=6pt, bottom=6pt
}

\setCJKmainfont{Hiragino Sans GB}
\setCJKsansfont{Heiti SC}

\begin{document}

% ============ 标题页 ============
\begin{center}
{\Huge\bfseries\textcolor{mainblue}{不等式与二次不等式}}\\[12pt]
{\Large\textcolor{gray}{初升高衔接 · 讲义三}}\\[8pt]
\rule{0.6\textwidth}{1pt}\\[8pt]
{\normalsize 适用：初升高衔接\quad |\quad 课时：2--3课时}
\end{center}
\vspace{10pt}

% ==============================================================
% 第一部分：知识点讲解
% ==============================================================
\section*{\textcolor{mainblue}{一、知识点讲解}}

% ---------- 1.1 不等式基本性质 ----------
\subsection*{1.1\quad 不等式基本性质}

\begin{knowledgebox}[不等式的三条基本性质]
设 $a, b, c$ 为实数，则：

\textbf{性质1（加法）：}若 $a > b$，则 $a + c > b + c$。

\quad 即：不等式两边同时加上（或减去）同一个数，不等号方向\textbf{不变}。

\textbf{性质2（乘正数）：}若 $a > b$，$c > 0$，则 $ac > bc$。

\quad 即：不等式两边同时乘以（或除以）同一个\textbf{正数}，不等号方向\textbf{不变}。

\textbf{性质3（乘负数）：}若 $a > b$，$c < 0$，则 $ac < bc$。

\quad 即：不等式两边同时乘以（或除以）同一个\textbf{负数}，不等号方向\textbf{改变}！
\end{knowledgebox}

\begin{warnbox}[易错提醒：性质3是最大坑！]
很多同学忘记"乘负数要变号"，这是不等式中最常见的错误。

\textbf{例子：}$-2x > 6$

错误做法：$x > -3$ \quad （直接除以$-2$，忘记变号）

正确做法：两边除以 $-2$，不等号反向，得 $x < -3$。

\medskip
\textbf{记忆口诀：}乘负变号，牢记不忘！
\end{warnbox}

% ---------- 1.2 一元一次不等式 ----------
\subsection*{1.2\quad 一元一次不等式回顾}

\begin{knowledgebox}[一元一次不等式的解法]
\textbf{步骤：}
\begin{enumerate}[label=(\arabic*)]
  \item 去分母（两边同乘最小公倍数，注意乘负数要变号）
  \item 去括号
  \item 移项（移项要变号）
  \item 合并同类项
  \item 系数化为1（除以负数要变号）
\end{enumerate}

\textbf{例子：}解不等式 $3x - 5 < 2x + 1$

移项：$3x - 2x < 1 + 5$，得 $x < 6$。

解集为 $\{x \mid x < 6\}$，用区间表示为 $(-\infty, 6)$。
\end{knowledgebox}

% ---------- 1.3 一元一次不等式组 ----------
\subsection*{1.3\quad 一元一次不等式组回顾}

\begin{knowledgebox}[一元一次不等式组的解法]
分别解出每个不等式的解集，再取\textbf{公共部分}。

\medskip
\textbf{口诀（四句话）：}

\begin{center}
\begin{tabular}{|c|c|c|}
\hline
\textbf{类型} & \textbf{条件} & \textbf{解集} \\
\hline
同大取大 & $\begin{cases} x > a \\ x > b \end{cases}$\quad ($a < b$) & $x > b$ \\
\hline
同小取小 & $\begin{cases} x < a \\ x < b \end{cases}$\quad ($a < b$) & $x < a$ \\
\hline
大小小大取中间 & $\begin{cases} x > a \\ x < b \end{cases}$\quad ($a < b$) & $a < x < b$ \\
\hline
大大小小无解 & $\begin{cases} x < a \\ x > b \end{cases}$\quad ($a < b$) & 无解 \\
\hline
\end{tabular}
\end{center}

\textbf{例子：}$\begin{cases} 2x - 1 > 3 \\ x + 2 < 7 \end{cases}$

解第一个：$x > 2$；解第二个：$x < 5$。

属于"大小小大取中间"，解集为 $2 < x < 5$，即 $(2, 5)$。
\end{knowledgebox}

% ---------- 1.4 一元二次不等式（重点！） ----------
\subsection*{1.4\quad 一元二次不等式（高中新内容，重点！）}

\begin{knowledgebox}[一元二次不等式的定义]
形如 $ax^2 + bx + c > 0$（或 $\geq 0$，$< 0$，$\leq 0$）的不等式，其中 $a \neq 0$，叫做\textbf{一元二次不等式}。

\medskip
\textbf{核心思想：}利用二次函数 $y = ax^2 + bx + c$ 的图象（抛物线）来求解。
\end{knowledgebox}

\begin{knowledgebox}[解法：图象法——"看图说话"]
\textbf{第一步：}化标准形式（$a > 0$，右边为 $0$）。

\textbf{第二步：}计算判别式 $\Delta = b^2 - 4ac$。

\textbf{第三步：}根据 $\Delta$ 的正负，画出抛物线的大致图象。

\textbf{第四步：}看抛物线在 $x$ 轴上方还是下方的部分，写出解集。

\medskip
设 $a > 0$，方程 $ax^2 + bx + c = 0$ 的两根为 $x_1, x_2$（$x_1 < x_2$）。
\end{knowledgebox}

\begin{knowledgebox}[$\Delta > 0$：抛物线与 $x$ 轴有两个交点]
\begin{center}
\begin{tikzpicture}
\begin{axis}[
  axis lines=middle, xlabel=$x$, ylabel=$y$,
  width=7cm, height=5cm,
  xmin=-1.5, xmax=4.5, ymin=-3, ymax=4,
  xtick={1, 3}, xticklabels={$x_1$, $x_2$},
  ytick=\empty,
  every axis x label/.style={at={(current axis.right of origin)},anchor=west},
  every axis y label/.style={at={(current axis.above origin)},anchor=south},
]
\addplot[domain=-0.5:4, samples=100, thick, mainblue] {1.5*(x-1)*(x-3)};
\draw[dashed, gray] (axis cs:1,0) -- (axis cs:1,-2.5);
\draw[dashed, gray] (axis cs:3,0) -- (axis cs:3,-2.5);
\node at (axis cs:0, 3.5) {$ax^2+bx+c=0$有两不等实根};
\end{axis}
\end{tikzpicture}
\end{center}

\begin{center}
\begin{tabular}{|c|c|}
\hline
\textbf{不等式} & \textbf{解集} \\
\hline
$ax^2+bx+c > 0$ & $\{x \mid x < x_1\text{ 或 }x > x_2\}$ \\
\hline
$ax^2+bx+c < 0$ & $\{x \mid x_1 < x < x_2\}$ \\
\hline
$ax^2+bx+c \geq 0$ & $\{x \mid x \leq x_1\text{ 或 }x \geq x_2\}$ \\
\hline
$ax^2+bx+c \leq 0$ & $\{x \mid x_1 \leq x \leq x_2\}$ \\
\hline
\end{tabular}
\end{center}
\end{knowledgebox}

\begin{knowledgebox}[$\Delta = 0$：抛物线与 $x$ 轴有唯一交点（相切）]
\begin{center}
\begin{tikzpicture}
\begin{axis}[
  axis lines=middle, xlabel=$x$, ylabel=$y$,
  width=7cm, height=5cm,
  xmin=-0.5, xmax=3.5, ymin=-0.5, ymax=4,
  xtick={2}, xticklabels={$x_0$},
  ytick=\empty,
  every axis x label/.style={at={(current axis.right of origin)},anchor=west},
  every axis y label/.style={at={(current axis.above origin)},anchor=south},
]
\addplot[domain=-0.2:3.2, samples=100, thick, mainblue] {2*(x-2)*(x-2)};
\draw[dashed, gray] (axis cs:2,0) -- (axis cs:2,-0.3);
\node at (axis cs:0.5, 3.5) {$ax^2+bx+c=0$有两相等实根};
\end{axis}
\end{tikzpicture}
\end{center}

设方程 $ax^2 + bx + c = 0$ 的根为 $x_0$（重根），则：

\begin{center}
\begin{tabular}{|c|c|}
\hline
\textbf{不等式} & \textbf{解集} \\
\hline
$ax^2+bx+c > 0$ & $\{x \mid x \neq x_0\}$（即全体实数去掉$x_0$） \\
\hline
$ax^2+bx+c < 0$ & 无解（$\varnothing$） \\
\hline
$ax^2+bx+c \geq 0$ & $\mathbb{R}$（全体实数） \\
\hline
$ax^2+bx+c \leq 0$ & $\{x_0\}$（只有一个点） \\
\hline
\end{tabular}
\end{center}
\end{knowledgebox}

\begin{knowledgebox}[$\Delta < 0$：抛物线与 $x$ 轴无交点]
\begin{center}
\begin{tikzpicture}
\begin{axis}[
  axis lines=middle, xlabel=$x$, ylabel=$y$,
  width=7cm, height=5cm,
  xmin=-1, xmax=4, ymin=-1, ymax=5,
  ytick=\empty, xtick=\empty,
  every axis x label/.style={at={(current axis.right of origin)},anchor=west},
  every axis y label/.style={at={(current axis.above origin)},anchor=south},
]
\addplot[domain=-0.5:3.5, samples=100, thick, mainblue] {(x-1.5)*(x-1.5)+1};
\node at (axis cs:1.5, 0.3) {$ax^2+bx+c=0$无实根};
\end{axis}
\end{tikzpicture}
\end{center}

\begin{center}
\begin{tabular}{|c|c|}
\hline
\textbf{不等式} & \textbf{解集} \\
\hline
$ax^2+bx+c > 0$ & $\mathbb{R}$（全体实数） \\
\hline
$ax^2+bx+c < 0$ & 无解（$\varnothing$） \\
\hline
$ax^2+bx+c \geq 0$ & $\mathbb{R}$ \\
\hline
$ax^2+bx+c \leq 0$ & 无解（$\varnothing$） \\
\hline
\end{tabular}
\end{center}
\end{knowledgebox}

\begin{tipbox}[核心口诀]
\textbf{"大于取两边，小于取中间"}（$\Delta > 0$，$a > 0$ 时）：

\begin{itemize}
  \item $ax^2+bx+c > 0$：解集取两根的\textbf{外侧}（两边）$x < x_1$ 或 $x > x_2$
  \item $ax^2+bx+c < 0$：解集取两根的\textbf{内侧}（中间）$x_1 < x < x_2$
\end{itemize}

\medskip
\textbf{注意：}如果 $a < 0$，先把两边乘以 $-1$（不等号变向），使 $a$ 变为正数再用口诀！
\end{tipbox}

% ---------- 1.5 绝对值不等式 ----------
\subsection*{1.5\quad 绝对值不等式}

\begin{knowledgebox}[绝对值不等式的两种基本类型]
设 $a > 0$，则：

\textbf{类型一：}$|x| < a$ \quad $\Longleftrightarrow$ \quad $-a < x < a$

\quad 含义：$x$ 到原点的距离\textbf{小于}$a$。

\textbf{类型二：}$|x| > a$ \quad $\Longleftrightarrow$ \quad $x < -a$ 或 $x > a$

\quad 含义：$x$ 到原点的距离\textbf{大于}$a$。

\medskip
\textbf{推广：}$|x| \leq a \Leftrightarrow -a \leq x \leq a$，\quad $|x| \geq a \Leftrightarrow x \leq -a$ 或 $x \geq a$。
\end{knowledgebox}

\begin{warnbox}[注意前提]
上面的等价关系\textbf{要求 $a > 0$}！

如果 $a \leq 0$，要单独讨论：
\begin{itemize}
  \item $|x| < a$ 且 $a \leq 0$：无解（绝对值不可能小于等于 $0$ 的数）
  \item $|x| > a$ 且 $a < 0$：解为全体实数（绝对值一定大于负数）
  \item $|x| > 0$：解为 $x \neq 0$
\end{itemize}
\end{warnbox}

% ---------- 1.6 分式不等式 ----------
\subsection*{1.6\quad 分式不等式（简单了解）}

\begin{knowledgebox}[分式不等式的基本解法]
\textbf{核心方法：}不要两边同乘分母（因为分母正负未知），而是\textbf{移项通分}，化为整式不等式。

\medskip
$\dfrac{f(x)}{g(x)} > 0$ \quad $\Longleftrightarrow$ \quad $f(x) \cdot g(x) > 0$（分子分母同号）

$\dfrac{f(x)}{g(x)} < 0$ \quad $\Longleftrightarrow$ \quad $f(x) \cdot g(x) < 0$（分子分母异号）

\medskip
\textbf{例子：}$\dfrac{x-1}{x+2} > 0$ \quad $\Longleftrightarrow$ \quad $(x-1)(x+2) > 0$

解得 $x < -2$ 或 $x > 1$。
\end{knowledgebox}

% ---------- 1.7 均值不等式预习 ----------
\subsection*{1.7\quad 均值不等式预习}

\begin{knowledgebox}[均值不等式（算术平均--几何平均不等式）]
若 $a \geq 0$，$b \geq 0$，则：
\[
\boxed{a + b \geq 2\sqrt{ab}}
\]
当且仅当 $a = b$ 时等号成立。

\medskip
\textbf{等价形式：}
\begin{itemize}
  \item 若 $a > 0$，则 $a + \dfrac{1}{a} \geq 2$
  \item 若 $ab > 0$，则 $\dfrac{a}{b} + \dfrac{b}{a} \geq 2$
\end{itemize}

\textbf{使用条件：}一正、二定、三相等
\begin{enumerate}[label=(\arabic*)]
  \item \textbf{一正：}$a, b$ 都是正数
  \item \textbf{二定：}$a + b$ 或 $ab$ 中至少有一个是定值
  \item \textbf{三相等：}等号能够取到（$a = b$）
\end{enumerate}
\end{knowledgebox}


% ==============================================================
% 第二部分：例题精讲
% ==============================================================
\section*{\textcolor{examorange}{二、例题精讲}}

% 例题1：基础不等式性质
\begin{examplebox}[【例题1】不等式基本性质的运用]
已知 $a > b$，判断下列结论是否正确：
\begin{enumerate}[label=(\arabic*)]
  \item $a + 3 > b + 3$
  \item $-2a > -2b$
  \item $a^2 > b^2$
\end{enumerate}
\end{examplebox}

\textbf{解：}
\begin{enumerate}[label=(\arabic*)]
  \item 正确。由性质1，两边加 $3$，不等号不变。
  \item 正确。由性质3，两边乘 $-2$，不等号变向。由 $a > b$ 得 $-2a < -2b$。

  \textbf{所以"$-2a > -2b$"是错误的！}

  \item 不一定正确。例如 $a = 1$，$b = -3$，满足 $a > b$，但 $a^2 = 1 < 9 = b^2$。
\end{enumerate}

\textbf{答：}(1) 正确；(2) 错误（乘负数应变号）；(3) 不一定正确。

\begin{tipbox}[方法总结]
判断不等式变形是否正确，关键是看：是否乘了负数（变号了吗？）、有没有特殊情况。举反例是否定一个命题的最有效方法。
\end{tipbox}

% 例题2：一元二次不等式（Δ > 0）
\begin{examplebox}[【例题2】一元二次不等式（$\Delta > 0$）]
解不等式：$x^2 - 5x + 6 > 0$
\end{examplebox}

\textbf{解：}

\textbf{第一步：}这里 $a = 1 > 0$，已经是标准形式。

\textbf{第二步：}求根。令 $x^2 - 5x + 6 = 0$，因式分解得 $(x-2)(x-3) = 0$。

解得 $x_1 = 2$，$x_2 = 3$。

\textbf{第三步：}因为 $a > 0$，抛物线开口向上，$\Delta = 25 - 24 = 1 > 0$，与 $x$ 轴有两个交点。

\textbf{第四步：}$> 0$ 取两边（大于取两边）。

\textbf{答：}解集为 $\{x \mid x < 2 \text{ 或 } x > 3\}$，即 $(-\infty, 2) \cup (3, +\infty)$。

% 例题3：一元二次不等式（Δ > 0，小于0）
\begin{examplebox}[【例题3】一元二次不等式（$\Delta > 0$，取中间）]
解不等式：$2x^2 + x - 6 \leq 0$
\end{examplebox}

\textbf{解：}

$a = 2 > 0$，已是标准形式。

令 $2x^2 + x - 6 = 0$，用求根公式：
\[
x = \frac{-1 \pm \sqrt{1 + 48}}{4} = \frac{-1 \pm 7}{4}
\]
得 $x_1 = -2$，$x_2 = \dfrac{3}{2}$。

$\Delta = 49 > 0$，开口向上，$\leq 0$ 取中间（含等号用闭区间）。

\textbf{答：}解集为 $\left\{x \mid -2 \leq x \leq \dfrac{3}{2}\right\}$，即 $\left[-2, \dfrac{3}{2}\right]$。

\begin{tipbox}[方法总结]
求根时如果不容易因式分解，就用求根公式 $x = \dfrac{-b \pm \sqrt{\Delta}}{2a}$。记住：$\leq 0$ 时两端点\textbf{取得到}（闭区间），$< 0$ 时两端点\textbf{取不到}（开区间）。
\end{tipbox}

% 例题4：Δ = 0 的情况
\begin{examplebox}[【例题4】一元二次不等式（$\Delta = 0$）]
解不等式：$x^2 - 6x + 9 > 0$
\end{examplebox}

\textbf{解：}

$a = 1 > 0$。

令 $x^2 - 6x + 9 = 0$，即 $(x - 3)^2 = 0$，解得 $x_0 = 3$。

$\Delta = 36 - 36 = 0$，抛物线与 $x$ 轴相切。

$(x-3)^2$ 永远 $\geq 0$，要使 $(x-3)^2 > 0$，只需 $x \neq 3$。

\textbf{答：}解集为 $\{x \mid x \neq 3\}$，即 $(-\infty, 3) \cup (3, +\infty)$。

% 例题5：Δ < 0 的情况
\begin{examplebox}[【例题5】一元二次不等式（$\Delta < 0$）]
解不等式：$x^2 + 2x + 5 > 0$
\end{examplebox}

\textbf{解：}

$a = 1 > 0$。

$\Delta = 4 - 20 = -16 < 0$，方程无实根，抛物线在 $x$ 轴\textbf{上方}。

所以 $x^2 + 2x + 5$ 恒大于 $0$。

\textbf{答：}解集为 $\mathbb{R}$（全体实数）。

\begin{tipbox}[三种情况速记表]
\begin{center}
\begin{tabular}{|c|c|c|}
\hline
$\Delta$ & $ax^2+bx+c > 0$ ($a>0$) & $ax^2+bx+c < 0$ ($a>0$) \\
\hline
$\Delta > 0$ & 取两边（两根外侧） & 取中间（两根之间） \\
\hline
$\Delta = 0$ & $x \neq x_0$ & 无解 \\
\hline
$\Delta < 0$ & 全体实数 & 无解 \\
\hline
\end{tabular}
\end{center}
\end{tipbox}

% 例题6：含参二次不等式
\begin{examplebox}[【例题6】含参二次不等式]
已知不等式 $x^2 - 2ax + a + 2 > 0$ 对一切实数 $x$ 恒成立，求 $a$ 的取值范围。
\end{examplebox}

\textbf{解：}

$a_0 = 1 > 0$，抛物线开口向上。

要使 $x^2 - 2ax + a + 2 > 0$ 恒成立，只需抛物线始终在 $x$ 轴上方，即：

$\Delta < 0$

$\Delta = (-2a)^2 - 4 \times 1 \times (a + 2) = 4a^2 - 4a - 8$

令 $4a^2 - 4a - 8 < 0$，即 $a^2 - a - 2 < 0$

因式分解：$(a - 2)(a + 1) < 0$

解得 $-1 < a < 2$。

\textbf{答：}$a$ 的取值范围为 $(-1, 2)$。

\begin{tipbox}[方法总结]
"恒大于 $0$"等价于"$\Delta < 0$"（开口向上时）。"恒小于 $0$"等价于"$\Delta < 0$ 且开口向下"。含参问题的核心是把参数当成常数处理，用判别式列不等式。
\end{tipbox}

% 例题7：绝对值不等式
\begin{examplebox}[【例题7】绝对值不等式]
解不等式：$|2x - 3| < 5$
\end{examplebox}

\textbf{解：}

设 $u = 2x - 3$，则 $|u| < 5$，即 $-5 < u < 5$。

代回：$-5 < 2x - 3 < 5$

各部分加 $3$：$-2 < 2x < 8$

各部分除以 $2$：$-1 < x < 4$

\textbf{答：}解集为 $(-1, 4)$。

\medskip
\textbf{另解：}$|2x - 3| < 5$ 直接等价于 $-5 < 2x - 3 < 5$，后续步骤相同。

% 例题8：不等式组与集合结合
\begin{examplebox}[【例题8】不等式组与集合结合（预习高中）]
设集合 $A = \{x \mid x^2 - 4x + 3 \leq 0\}$，$B = \{x \mid |x - 2| < 2\}$，求 $A \cap B$。
\end{examplebox}

\textbf{解：}

先求集合 $A$：解 $x^2 - 4x + 3 \leq 0$

$(x - 1)(x - 3) \leq 0$，解得 $1 \leq x \leq 3$。

所以 $A = [1, 3]$。

\medskip
再求集合 $B$：解 $|x - 2| < 2$

$-2 < x - 2 < 2$，解得 $0 < x < 4$。

所以 $B = (0, 4)$。

\medskip
求交集 $A \cap B$：取 $[1, 3]$ 和 $(0, 4)$ 的公共部分。

\textbf{答：}$A \cap B = [1, 3]$（因为 $[1, 3] \subset (0, 4)$）。

\begin{tipbox}[方法总结]
不等式与集合结合是高中常见题型。解法：分别解出各集合中的不等式，再用数轴或区间取交集、并集。记住：交集取公共部分，并集取合在一起。
\end{tipbox}

% 例题9：图象法综合
\begin{examplebox}[【例题9】图象法解不等式]
已知二次函数 $y = ax^2 + bx + c$ 的图象如图所示，写出 $ax^2 + bx + c > 0$ 和 $ax^2 + bx + c \leq 0$ 的解集。

\begin{center}
\begin{tikzpicture}
\begin{axis}[
  axis lines=middle, xlabel=$x$, ylabel=$y$,
  width=8cm, height=5cm,
  xmin=-3, xmax=5, ymin=-4, ymax=4,
  xtick={-1, 3}, xticklabels={$-1$, $3$},
  ytick=\empty,
  every axis x label/.style={at={(current axis.right of origin)},anchor=west},
  every axis y label/.style={at={(current axis.above origin)},anchor=south},
]
\addplot[domain=-2.5:4.5, samples=100, thick, mainblue] {-1*(x+1)*(x-3)};
\draw[dashed, gray] (axis cs:-1,0) -- (axis cs:-1,-3);
\draw[dashed, gray] (axis cs:3,0) -- (axis cs:3,-3);
\end{axis}
\end{tikzpicture}
\end{center}
\end{examplebox}

\textbf{解：}

从图象可知：
\begin{itemize}
  \item 抛物线\textbf{开口向下}，所以 $a < 0$。
  \item 与 $x$ 轴交于 $x = -1$ 和 $x = 3$。
\end{itemize}

\textbf{（1）}$ax^2 + bx + c > 0$：图象在 $x$ 轴\textbf{上方}的部分。

\quad 从图上看，当 $-1 < x < 3$ 时，$y > 0$。

\quad 解集为 $(-1, 3)$。

\textbf{（2）}$ax^2 + bx + c \leq 0$：图象在 $x$ 轴\textbf{下方或在轴上}的部分。

\quad 从图上看，当 $x \leq -1$ 或 $x \geq 3$ 时，$y \leq 0$。

\quad 解集为 $(-\infty, -1] \cup [3, +\infty)$。

\begin{tipbox}[方法总结]
当 $a < 0$ 时，口诀刚好\textbf{相反}："大于取中间，小于取两边"。也可以先把不等式两边乘以 $-1$（变号），使 $a$ 变为正数，再用正常口诀。
\end{tipbox}


% ==============================================================
% 第三部分：应用题
% ==============================================================
\section*{\textcolor{apppurple}{三、应用题}}

\begin{appbox}[【应用1】商品定价问题]
某商店以每件 $20$ 元的成本购进一种商品。据市场调查：当售价为每件 $x$ 元（$x \geq 20$）时，每天可卖出 $(100 - 2x)$ 件。要使每天的利润不低于 $450$ 元，售价应定在什么范围？
\end{appbox}

\textbf{解：}

每件利润为 $(x - 20)$ 元，每天卖出 $(100 - 2x)$ 件。

每天总利润 $= (x - 20)(100 - 2x)$

令 $(x - 20)(100 - 2x) \geq 450$

展开：$100x - 2x^2 - 2000 + 40x \geq 450$

$-2x^2 + 140x - 2000 \geq 450$

$-2x^2 + 140x - 2450 \geq 0$

两边除以 $-2$（变号）：$x^2 - 70x + 1225 \leq 0$

$(x - 35)^2 \leq 0$

因为 $(x - 35)^2 \geq 0$，所以 $(x - 35)^2 \leq 0$ 只在 $x = 35$ 时成立。

\textbf{答：}售价应定为 $35$ 元/件。

\textbf{（验证：}$x = 35$ 时，卖出 $100 - 70 = 30$ 件，每件利润 $15$ 元，总利润 $= 30 \times 15 = 450$ 元。）

\begin{appbox}[【应用2】围栏面积问题]
用总长为 $40$ 米的篱笆围成一个一面靠墙的矩形花园（墙长足够长）。要使花园面积不小于 $180$ 平方米，求平行于墙的一边长度的范围。（靠墙的一边不需篱笆）
\end{appbox}

\textbf{解：}

设平行于墙的一边长度为 $x$ 米，则另一边长度为 $\dfrac{40 - x}{2}$ 米。

由题意 $x > 0$ 且 $\dfrac{40 - x}{2} > 0$，得 $0 < x < 40$。

面积 $S = x \cdot \dfrac{40 - x}{2} = \dfrac{40x - x^2}{2}$

令 $S \geq 180$：

$\dfrac{40x - x^2}{2} \geq 180$

$40x - x^2 \geq 360$

$-x^2 + 40x - 360 \geq 0$

$x^2 - 40x + 360 \leq 0$

用求根公式：$x = \dfrac{40 \pm \sqrt{1600 - 1440}}{2} = \dfrac{40 \pm \sqrt{160}}{2} = \dfrac{40 \pm 4\sqrt{10}}{2} = 20 \pm 2\sqrt{10}$

$\sqrt{10} \approx 3.16$，$2\sqrt{10} \approx 6.32$。

$x_1 \approx 13.68$，$x_2 \approx 26.32$。

$\leq 0$ 取中间：$20 - 2\sqrt{10} \leq x \leq 20 + 2\sqrt{10}$。

\textbf{答：}平行于墙的一边长度范围为 $[20 - 2\sqrt{10},\ 20 + 2\sqrt{10}]$（约 $[13.7, 26.3]$ 米）。

\begin{appbox}[【应用3】利润不低于目标]
某工厂生产一种产品，固定成本为 $1000$ 元，每生产一件产品需再投入 $20$ 元。若产品售价为每件 $45$ 元，且全部售出。
\begin{enumerate}[label=(\arabic*)]
  \item 写出利润 $P$（元）与产量 $x$（件）的函数关系。
  \item 若要盈利（即 $P > 0$），至少需要生产多少件？
  \item 若要利润不低于 $1500$ 元，产量应为多少？
\end{enumerate}
\end{appbox}

\textbf{解：}

\textbf{（1）}$P = 45x - (1000 + 20x) = 25x - 1000$

\textbf{（2）}令 $P > 0$：$25x - 1000 > 0$，解得 $x > 40$。

因为 $x$ 为整数，所以至少生产 $41$ 件。

\textbf{（3）}令 $P \geq 1500$：$25x - 1000 \geq 1500$

$25x \geq 2500$，解得 $x \geq 100$。

\textbf{答：}(1) $P = 25x - 1000$；(2) 至少 $41$ 件；(3) 产量 $\geq 100$ 件。


% ==============================================================
% 第四部分：课堂练习
% ==============================================================
\section*{\textcolor{tipgreen}{四、课堂练习}}

\begin{enumerate}
  \item 解不等式：$3x + 5 > 2x - 1$

  \item 解不等式组：$\begin{cases} 3x - 1 > 5 \\ 7 - 2x > 1 \end{cases}$

  \item 解不等式：$x^2 - 3x - 4 > 0$

  \item 解不等式：$-x^2 + 4x - 3 \geq 0$

  \item 解不等式：$2x^2 - 4x + 3 < 0$

  \item 解不等式：$|3x + 1| \leq 7$

  \item 解不等式：$\dfrac{x + 1}{x - 3} \leq 0$

  \item 已知关于 $x$ 的不等式 $x^2 - 2mx + m + 2 > 0$ 对任意实数 $x$ 恒成立，求 $m$ 的取值范围。
\end{enumerate}

\vspace{20pt}
\begin{knowledgebox}[参考答案]

\textbf{1.} $3x - 2x > -1 - 5$，得 $x > -6$。解集：$(-6, +\infty)$。

\medskip
\textbf{2.} 解第一个：$3x > 6$，$x > 2$。解第二个：$-2x > -6$，$x < 3$。

大小小大取中间，解集：$(2, 3)$。

\medskip
\textbf{3.} $(x - 4)(x + 1) > 0$，$a > 0$，大于取两边。

$x_1 = -1$，$x_2 = 4$。解集：$(-\infty, -1) \cup (4, +\infty)$。

\medskip
\textbf{4.} 两边乘 $-1$（变号）：$x^2 - 4x + 3 \leq 0$

$(x - 1)(x - 3) \leq 0$，解集：$[1, 3]$。

\medskip
\textbf{5.} $\Delta = 16 - 24 = -8 < 0$，$a = 2 > 0$，抛物线全在 $x$ 轴上方。

$2x^2 - 4x + 3$ 恒大于 $0$，解集：无解（$\varnothing$）。

\medskip
\textbf{6.} $|3x + 1| \leq 7$ $\Leftrightarrow$ $-7 \leq 3x + 1 \leq 7$

$-8 \leq 3x \leq 6$，得 $-\dfrac{8}{3} \leq x \leq 2$。解集：$\left[-\dfrac{8}{3},\ 2\right]$。

\medskip
\textbf{7.} $\dfrac{x + 1}{x - 3} \leq 0$ $\Leftrightarrow$ $(x + 1)(x - 3) \leq 0$ 且 $x \neq 3$。

由 $(x+1)(x-3) \leq 0$ 得 $-1 \leq x \leq 3$，再去掉 $x = 3$。

解集：$[-1, 3)$。

\medskip
\textbf{8.} 恒成立需 $\Delta < 0$。

$\Delta = 4m^2 - 4(m + 2) = 4m^2 - 4m - 8$

令 $4m^2 - 4m - 8 < 0$，即 $m^2 - m - 2 < 0$

$(m - 2)(m + 1) < 0$，解得 $-1 < m < 2$。

$m$ 的取值范围：$(-1, 2)$。
\end{knowledgebox}

\vfill
\begin{center}
\small\textcolor{gray}{本讲义仅供教学参考，例题可根据学生水平调整难度。}
\end{center}

\end{document}
